\begin{notechapter}
\chapter{Word Embeddings}\label{ch:word-embeddings}
\chaptersubtitle{Word2vec, GloVe and fastText: vectors learned from the company words keep}

\begin{chapterintro}
In the neural bigram model, each word received a vector, and words with similar
continuations received similar vectors (\cref{ch:bigram-neural}). Word embeddings make this
the goal: a prediction task over a large corpus is set up only to learn one vector per word,
and the vectors are then used as features by other models. This chapter states the
models precisely, with their two sets of vectors, their objectives and their gradients,
and is careful about two claims that are often made loosely: what the vector arithmetic of
analogies shows, and which tools actually provide static word vectors. You will learn
\begin{itemize}[nosep]
  \item the distributional hypothesis and the skip-gram and CBOW models of word2vec, with
        their input and output vectors and the cost of their softmax;
  \item the GloVe objective, a weighted least-squares fit of log co-occurrence counts;
  \item fastText's character $n$-grams and how they handle unseen words;
  \item how analogies are evaluated, and why the query words must be excluded;
  \item static versus contextual embeddings, and how embeddings are stored and loaded.
\end{itemize}
\prerequisites{\cref{ch:bigram-neural}; the gradient of the softmax cross-entropy
(\cref{bigram-models:sec:gradient}).}
\end{chapterintro}

\section{The distributional hypothesis}\label{word-embeddings:sec:distributional}

``You shall know a word by the company it keeps'' (Firth, 1957). The \emph{distributional
hypothesis} (Harris, 1954) states that words occurring in similar contexts tend to have
similar meanings: ``coffee'' and ``tea'' both appear near ``cup'', ``drink'' and ``hot''. A word
embedding is a map $w\mapsto v_w\in\R^d$, with $d$ of a few hundred, learned so that words with similar
contexts get nearby vectors. Similarity is measured by the cosine,
$\cos(v_a, v_b) = v_a^\top v_b/(\norm{v_a}\,\norm{v_b})$.

The models below learn the vectors through a \emph{pretext task}: predicting words from their
neighbours. The predictions themselves are discarded after training; only the vectors are
kept. This is the idea of self-supervised learning (\cref{part:contrastive}): the
supervision comes from the data itself, here from the order of words in raw text.

\section{Word2vec}\label{word-embeddings:sec:word2vec}

Word2vec (Mikolov, Chen, Corrado and Dean, \emph{Efficient Estimation of Word Representations
in Vector Space}, 2013) has two architectures, both with two matrices: input vectors $v_w$ (rows
of $V_{\mathrm{in}}\in\R^{|\mathcal V|\times d}$) and output vectors $u_w$ (rows of $U\in\R^{|\mathcal V|\times d}$), where $\mathcal V$ is the vocabulary. A text $w_1,\dots,w_N$ and a window
size $m$ define the training pairs: each \emph{centre} word $w_t$ with each \emph{context} word $w_{t+k}$, $1\leq\abs k\leq m$.

\paragraph{Skip-gram} predicts each context word from the centre word:
\[
  p(c\mid w) = \frac{\exp(u_c^\top v_w)}{\sum_{c'\in\mathcal V}\exp(u_{c'}^\top v_w)},\qquad
  \mathcal L = -\sum_{t=1}^N\sum_{1\leq\abs k\leq m}\log p(w_{t+k}\mid w_t) .
\]
This is the neural bigram model of \cref{ch:bigram-neural} without the hidden nonlinearity and
without biases, with the next word replaced by any word of the window: the logits are the
dot products $z = Uv_w$, and the model is \emph{log-bilinear}. For one pair $(w, c)$, the gradients follow from
$\nabla_z\ell = p - e_c$ (\cref{bigram-models:sec:gradient}):
\[
  \nabla_{v_w}\ell = U^\top(p - e_c) = \sum_{c'}p(c'\mid w)\,u_{c'} - u_c,\qquad
  \nabla_{u_{c'}}\ell = \bigl(p(c'\mid w) - [c' = c]\bigr)v_w .
\]
The centre vector moves towards the output vector of the observed context and away from
the model's average output vector; the observed context's output vector moves towards
$v_w$, and every other output vector moves away from it.

\paragraph{CBOW} (continuous bag of words) predicts the centre word from the average of the
input vectors of its context, $\bar v_t = \frac1{2m}\sum_{1\leq\abs k\leq m}v_{w_{t+k}}$:
$p(w_t\mid\text{context}) = \softmax(U\bar v_t)_{w_t}$. Averaging the vectors is the same as multiplying $V_{\mathrm{in}}^\top$ by the average of the
one-hot codes, since the lookup is linear; the order of the context words is lost, hence ``bag
of words''.

\begin{figure}[ht]
\centering
\begin{tikzpicture}[
  box/.style={draw, rounded corners=2pt, minimum height=8mm, inner sep=3pt, font=\small, align=center},
  ar/.style={-{Stealth[length=1.8mm]}, thick},
  lab/.style={font=\scriptsize, inner sep=1pt}]
  \node[box, fill=bookaccent!15] (w) at (0,0) {centre word $w$\\[-1pt]{\scriptsize index}};
  \node[box, fill=bookexample!15] (v) at (3.7,0) {$v_w\in\R^d$\\[-1pt]{\scriptsize row of $V_{\mathrm{in}}$}};
  \node[box, fill=bookkey!12] (z) at (7.5,0) {$z = Uv_w$\\[-1pt]{\scriptsize $|\mathcal V|$ dot products}};
  \node[box, fill=gray!10] (p) at (11.2,0) {$p(\cdot\mid w) = \softmax(z)$};
  \draw[ar] (w) -- node[lab, above]{lookup} (v);
  \draw[ar] (v) -- node[lab, above]{$U$} (z);
  \draw[ar] (z) -- (p);
  \node[box, fill=bookexample!15, dashed] (c) at (11.2,-1.4) {target: each context word $w_{t+k}$};
  \draw[ar, dashed, gray] (c) -- (p);
\end{tikzpicture}
\caption{Skip-gram. The input vectors are the embeddings that are usually kept; the softmax over the
whole vocabulary is the expensive part.}
\label{word-embeddings:fig:skipgram}
\end{figure}

\paragraph{The cost of the softmax.} Each pair costs $O(d)$ for the lookup and $O(d\,|\mathcal V|)$ for the logits and the
normalization, with $|\mathcal V|$ in the hundreds of thousands. Word2vec never computes the full softmax.
It replaces it either by a hierarchical softmax (\cref{ch:hierarchical-softmax}) or by negative
sampling (\cref{ch:nce-negative-sampling}), both $O(d\log|\mathcal V|)$ or $O(dk)$ per pair. Two further details of the
implementation matter in practice: frequent words are randomly discarded (subsampling,
each occurrence of $w$ dropped with probability $1 - \sqrt{s/f(w)}$, with $f(w)$ the frequency of $w$ and $s$ a threshold
around $10^{-5}$), and the window size is drawn at random between $1$ and $m$ for each centre word, which
gives nearby words more weight (Mikolov, Sutskever, Chen, Corrado and Dean,
\emph{Distributed Representations of Words and Phrases and their Compositionality}, 2013).

\paragraph{Which vectors?} The model has two vectors per word, and they play different roles:
$v_w$ describes $w$ as a centre, $u_w$ as a context. The input vectors are the ones usually published and
used; some methods use $v_w + u_w$. The two are not interchangeable: words that are \emph{similar} have
similar input vectors, while a large $u_c^\top v_w$ means that $c$ is likely \emph{near} $w$, which holds for
``coffee'' and ``cup'' as much as for ``coffee'' and ``tea''.

\section{GloVe}\label{word-embeddings:sec:glove}

GloVe (Pennington, Socher and Manning, \emph{GloVe: Global Vectors for Word Representation},
2014) starts from the co-occurrence counts. Let $X_{ij}$ be the number of times word $j$ occurs in the
window of word $i$ over the whole corpus (each occurrence weighted by $1/\text{distance}$). GloVe
minimizes
\[
  J = \sum_{i,j:\,X_{ij}>0}f(X_{ij})\,\bigl(w_i^\top\tilde w_j + b_i + \tilde b_j - \log X_{ij}\bigr)^2,\qquad
  f(x) = \begin{cases}(x/x_{\max})^{3/4} & x < x_{\max},\\ 1 & x\geq x_{\max},\end{cases}
\]
with $x_{\max} = 100$ in the paper. It is a weighted least-squares fit of the log-counts by a
low-rank matrix plus biases. The sum runs over the non-zero counts only, since $\log0$ is undefined,
which also makes each pass cost the number of distinct co-occurring pairs rather than
$|\mathcal V|^2$. The weight $f$ does two things: it gives little weight to rare, noisy co-occurrences, and it
caps the weight of very frequent pairs (``of the'') at $1$, so that they do not dominate. The
published vectors are $w_i + \tilde w_i$. Unlike word2vec, GloVe has no softmax, so neither hierarchical
softmax nor negative sampling plays any role in it.

\section{fastText: subword vectors}\label{word-embeddings:sec:fasttext}

Word2vec and GloVe learn one vector per vocabulary entry; a word absent from the training
corpus has none, and rare words have poor ones. FastText (Bojanowski, Grave, Joulin and
Mikolov, \emph{Enriching Word Vectors with Subword Information}, 2017) represents a word by
its character $n$-grams. The word is wrapped in boundary symbols, \texttt{<apple>}, and all its $n$-grams
with $3\leq n\leq6$ are extracted, plus the whole word itself as a special sequence: for $n = 3$,
\texttt{<ap}, \texttt{app}, \texttt{ppl}, \texttt{ple}, \texttt{le>}. The vector of the word is the sum of the
vectors of its $n$-grams, and these are trained with the skip-gram objective and negative
sampling. The boundary symbols distinguish the trigram \texttt{her} inside ``where'' from the
word \texttt{<her>}. An unseen word still has $n$-grams, hence a vector; words sharing a stem or a
suffix share parameters, which helps for rare words and morphologically rich languages.
To bound memory, the $n$-grams are hashed into a fixed number of buckets (two million by
default), so that two $n$-grams may share a vector.

\section{Analogies and similarity}\label{word-embeddings:sec:analogies}

Word2vec's vectors became famous for analogies: the vector $v_{\text{king}} - v_{\text{man}} + v_{\text{woman}}$ is close to $v_{\text{queen}}$,
and $v_{\text{Paris}} - v_{\text{France}} + v_{\text{Italy}}$ to $v_{\text{Rome}}$. The standard evaluation answers ``$a$ is to $b$ as $c$ is to ?'' by
\[
  \hat d = \operatorname*{arg\,max}_{x\in\mathcal V\setminus\{a,b,c\}}\cos\bigl(v_x,\ v_b - v_a + v_c\bigr),
\]
and the set $\mathcal V\setminus\{a,b,c\}$ matters.

\begin{pitfall}[The query words are excluded]
The offset $v_b - v_a$ between ``man'' and ``king'' is often small compared with the vectors
themselves, so the nearest neighbour of $v_{\text{king}} - v_{\text{man}} + v_{\text{woman}}$ is very often ``king'' itself, or
``woman''. The evaluation of Mikolov et al.\ excludes the three query words, and the
reported accuracies depend on it (Nissim, van Noord and van der Goot, \emph{Fair is Better
than Sensational: Man is to Doctor as Woman is to Doctor}, 2020). An analogy shows that the
expected answer is the nearest \emph{other} word, which is weaker than ``the vector arithmetic
produces queen''.
\end{pitfall}

\begin{example}[Excluding the query words]\label{word-embeddings:ex:analogy}
Take the three-dimensional vectors man $= (1, 0.3, 0.2)$, woman $= (1, -0.3, 0.2)$, king $= (3, 0.4, 1)$,
queen $= (2.6, -0.3, 1.6)$ and prince $= (2.8, 0.5, 0.6)$. Then $q = v_{\text{king}} - v_{\text{man}} + v_{\text{woman}} = (3, -0.2, 1)$, and the cosines
with $q$ are $0.9822$ for king, $0.9733$ for queen, $0.9679$ for woman, $0.9663$ for prince and $0.9322$ for man. Over
the whole vocabulary the answer is ``king''; with the query words excluded it is ``queen''.
\end{example}

\section{Static and contextual embeddings}\label{word-embeddings:sec:contextual}

Word2vec, GloVe and fastText are \emph{static}: one vector per word type, whatever the sentence.
``Bank'' gets a single vector, an average of its financial and riverside uses. \emph{Contextual}
embeddings are computed for each occurrence by a deep network that reads the whole
sentence. ELMo (Peters et al., \emph{Deep Contextualized Word Representations}, 2018) combines
the hidden states of a two-layer bidirectional LSTM language model (\cref{sequence-models:sec:bidirectional})
with learned task-specific weights. BERT (Devlin et al., 2019) uses the hidden states of a
transformer encoder trained to predict masked words, over subword units rather than words.
The static models are shallow: their vectors are rows of a parameter matrix, not activations
of a deep network; the contextual ones are activations, and depend on the sentence.

\section{Storing and loading embeddings}\label{word-embeddings:sec:implementation}

An embedding is stored as two objects: a mapping from words to indices, and a matrix with
one row per index.
\begin{lstlisting}[language=Python]
index_mapping = {"king": 0, "queen": 1, "man": 2, "woman": 3}
embedding_matrix = np.array([[0.7, 0.8, 0.9],    # row 0: "king"
                             [0.6, 0.7, 0.8],    # row 1: "queen"
                             [0.5, 0.6, 0.7],    # row 2: "man"
                             [0.4, 0.5, 0.6]])   # row 3: "woman"
vector = embedding_matrix[index_mapping["king"]] # array([0.7, 0.8, 0.9])
\end{lstlisting}
Storing the vectors in one matrix rather than in a dictionary of separate arrays keeps them
contiguous in memory, lets a whole batch be looked up with one indexing operation (an
array of indices in place of one index), and makes the table a single trainable parameter.
In PyTorch it is \texttt{nn.Embedding(num\_words, d)}: the backward pass of a lookup adds the
gradient of each looked-up vector to its row and leaves the other rows untouched, so the
gradient is sparse, as in \cref{bigram-models:sec:gradient}. Pretrained vectors initialize it with
\texttt{nn.Embedding.from\_pretrained(tensor, freeze=True)}.

With \texttt{gensim}, pretrained word2vec or fastText vectors are loaded as \texttt{KeyedVectors},
which implement the analogy query with the exclusion of the query words:
\begin{lstlisting}[language=Python]
from gensim.models import KeyedVectors
kv = KeyedVectors.load_word2vec_format("GoogleNews-vectors-negative300.bin", binary=True)
kv["king"].shape                                   # (300,)
kv.most_similar(positive=["king", "woman"], negative=["man"], topn=3)
\end{lstlisting}

\begin{pitfall}[spaCy's small models have no word vectors]
In spaCy, \texttt{nlp(word).vector} returns a vector for any model, but the small English
pipeline \texttt{en\_core\_web\_sm} ships without static word vectors: the vectors it returns
are context-sensitive features of its tagger, not word embeddings, and similarities computed
from them are unreliable. Static vectors come with the medium and large pipelines
(\texttt{en\_core\_web\_md}, \texttt{en\_core\_web\_lg}); \texttt{token.has\_vector} tells which case applies.
\end{pitfall}

\begin{chaptersummary}
\begin{itemize}
  \item Embeddings map words to vectors so that words with similar contexts are close; they
        are learned through a pretext task on raw text.
  \item Skip-gram predicts context words from the centre word with $p(c\mid w)\propto\exp(u_c^\top v_w)$; CBOW predicts the
        centre from the average context vector. Each word has an input and an output vector;
        the full softmax costs $O(d|\mathcal V|)$ per pair and is replaced in practice.
  \item GloVe fits $w_i^\top\tilde w_j + b_i + \tilde b_j\approx\log X_{ij}$ by weighted least squares over the non-zero co-occurrence counts.
  \item FastText sums the vectors of character $n$-grams ($3\leq n\leq6$, with boundary symbols), which gives
        vectors to unseen words.
  \item Analogy accuracy is measured with the query words excluded; without the exclusion the
        answer is often a query word.
  \item Static embeddings give one vector per word; contextual ones (ELMo, BERT) one per
        occurrence.
\end{itemize}
\end{chaptersummary}

\section*{Exercises}
\addcontentsline{toc}{section}{Exercises}

\begin{exercise}\label{word-embeddings:exr:1}
List the skip-gram training pairs of the sentence \emph{the cat sat on the mat} with $m = 1$. How many
are there with $m = 2$?
\solutionlink{word-embeddings:exr:1}
\end{exercise}

\begin{exercise}\label{word-embeddings:exr:2}
Show that the CBOW input $\bar v_t$ equals $V_{\mathrm{in}}^\top\bar x_t$, where $\bar x_t$ is the average of the one-hot codes of the context
words. What happens if a context word appears twice in the window?
\solutionlink{word-embeddings:exr:2}
\end{exercise}

\begin{exercise}\label{word-embeddings:exr:3}
Compute the GloVe weights $f(X_{ij})$ for $X_{ij} = 1, 10, 100, 1000$. Why is the pair with $X_{ij} = 0$ not in the sum,
and why would it be a problem to include it with $\log(X_{ij} + 1)$ and weight $1$?
\solutionlink{word-embeddings:exr:3}
\end{exercise}

\begin{exercise}\label{word-embeddings:exr:4}
List the fastText $n$-grams of the word \emph{where} for $n = 3$ and $n = 4$. Which of them does it share with the word
\emph{her}?
\solutionlink{word-embeddings:exr:4}
\end{exercise}

\begin{exercise}\label{word-embeddings:exr:5}
Let $\delta = v_c - v_a$, and assume that no other word has a vector with the direction of $v_b$. Show that, as
$\norm\delta\to0$, the word whose vector has the largest cosine with $v_b + \delta$ is eventually $b$. Why does this justify
excluding $b$ from the search?
\solutionlink{word-embeddings:exr:5}
\end{exercise}
\end{notechapter}
